\chapter{Судалгааны арга зүй}

\section{Системийн зохиомж}

Энэхүү бүлэгт судалгааны ерөнхий зорилго, дэвшүүлсэн таамаглал, хариулах судалгааны асуултууд болон системийн хэрэгжилтийн аргачлалыг тодорхойлно. Судалгааны гол зорилго нь камер болон хиймэл оюун ухаанд суурилсан ирц бүртгэлийн системийн архитектур боловсруулах, бодит орчинд туршин гүйцэтгэлийг тоон үзүүлэлтээр үнэлэхэд оршино.

Судалгааны үндсэн асуултууд нь:
\begin{itemize}
    \item Бодит анги танхимын нөхцөлд системийн Accuracy, Recall, FAR зэрэг үзүүлэлтүүд хүлээн зөвшөөрөгдөх түвшинд хүрэх үү?
    \item Аль алгоритм болон архитектурын шийдэл нь хамгийн сайн гүйцэтгэл үзүүлэх вэ? 
    \item Системийн гүйцэтгэлийн хугацааны үзүүлэлтүүд (нэг зураг боловсруулах хугацаа, нэг нүүр таних хугацаа) бодит хэрэглээнд тохиромжтой байх уу?
    \item Системийн практик хэрэглээнд тулгарч буй хязгаарлалтууд болон сайжруулалтын боломжууд юу вэ?
\end{itemize}

\subsection{Системийн архитектур}

Систем нь дүрс хүлээн авах модуль, нүүр илрүүлэх хэсэг, эмбеддинг үүсгэх гүн сургалтын загвар, kNN, Cosine Similarity ашигласан hybrid ангилагч, ирц баталгаажуулах механизм, өгөгдөл хадгалах сан гэсэн үндсэн бүрэлдэхүүнүүдээс тогтоно. Өгөгдлийн урсгал нь зураг авах үеэс эхлэн ирцийн CSV тайлан үүсэх хүртэл шаталсан боловсруулалтаар дамжина (Зураг~\ref{fig:system-arch}).

\begin{figure}[H]
  \centering
  \includegraphics[width=\textwidth]{diagrams/SystemDiagram.png}
  \caption{Системийн архитектурын ерөнхий схем}
  \label{fig:system-arch}
\end{figure}
\FloatBarrier

\subsection{Ирц бүртгэх процессын дараалал}

Систем нь бодит хугацааны видео урсгалыг шууд боловсруулахын оронд тогтмол давтамжтайгаар зураг авч боловсруулна. Энэ арга нь тооцооллын ачааллыг бууруулж, шийдвэрийг олон ажиглалтад тулгуурлан баталгаажуулах боломж олгоно (Зураг~\ref{fig:sequence-diagram}).

Нэг оюутан тухайн хичээлийн хугацаанд $H$-аас дээш удаа тогтвортой танигдсан тохиолдолд ирцэд орсон гэж үзнэ.

\begin{figure}[H]
  \centering
  \includegraphics[width=\textwidth,height=0.78\textheight,keepaspectratio]{diagrams/SequenceDiagram.png}
  \caption{Ирц бүртгэх процессын дарааллын диаграм}
  \label{fig:sequence-diagram}
\end{figure}
\FloatBarrier

\subsection{Нүүр илрүүлэлтийн шат}

Камераас хүлээн авсан RGB зураг бүр дээр дараах шаталсан боловсруулалт хийгдэнэ:

\begin{enumerate}
  \item Зураг хүлээн авах.
  \item Нүүрний хүрээ (bounding box) болон итгэлцүүрийн оноо(confidence score) тооцоолох.
  \item Гол тулгуур цэгүүд (нүд, хамар, ам) илрүүлэх.
  \item Геометрийн нормчилол (эргүүлэлт, шилжилт).
  \item Тайрах ба стандарт хэмжээ (\(160\times160\)) рүү хөрвүүлэх.
  \item Чанарын шалгалт (хэмжээ, бүдгэрэл, илрүүлэлтийн магадлал).
  \item Аннотацтай(хэдэн хувьтай таньсан эсэх) хадгалалт.
\end{enumerate}

Зөвшөөрөгдсөн нүүрнүүд дараагийн эмбеддинг үүсгэх шат руу шилжинэ.

\begin{figure}[H]
  \centering
  \includegraphics[width=0.95\textwidth]{ClassToCropSequence_IMG6373.jpg}
  \caption{Нүүр илрүүлэлтээс эмбеддингт бэлтгэх шаталсан жишээ}
  \label{fig:process-class-close-landmark-crop}
\end{figure}

\begin{figure}[H]
  \centering
  \includegraphics[width=\textwidth]{EmbeddingProcess.jpg}
  \caption{Тайрсан зургаас эмбеддинг вектор үүсэх CNN загварын дүрслэл}
  \label{fig:crop-to-embedding}
\end{figure}

\subsection{Hybrid царай таних арга}

Системийн таних механизм нь cosine similarity болон k-NN үнэлгээг хослуулсан Hybrid арга юм. Сурагчдын царайны эмбеддинг векторын сан (галлерей)-тай харьцуулж дараах шатлал хэрэгжинэ:

\begin{enumerate}
  \item Эмбеддинг векторыг L2 нормчлох.
  \item $k$-ойр хөршүүдийг сонгож, ангилал бүрийн нийлбэр төстэй байдлыг тооцох.
  \item Төв эмбеддингтэй cosine төстэй байдлыг тооцох.
  \item Нэгдсэн оноог дараах байдлаар тодорхойлох:
  \[
  S_{hyb} = \alpha S_{knn} + (1-\alpha) S_{cos}
  \]
  \item Итгэлцүүрийн босгоор \texttt{unknown} ангиллыг шүүх.
  \item Чанарын шалгалт хийж, шаардлага хангахгүй тохиолдлыг хасах.
\end{enumerate}

Hybrid буюу kNN болон cosine similarity хосолсон арга нь олон ажиглалтад тогтвортой гүйцэтгэл үзүүлсэн тул үндсэн таних механизм болгон сонгосон.

\subsection{Ирц баталгаажуулах механизм}

Ирц баталгаажуулалтыг олон давтамжит ажиглалтад тулгуурлан гүйцэтгэнэ:

\begin{enumerate}
  \item $\Delta t$ хугацаанд нэг зураг авах.
  \item Танилтын оноо $\tau$-аас дээш бол нэг \textit{hit} бүртгэх.
  \item Давхардлаас сэргийлэх $g$ хугацааны шүүлт.
  \item $h_i \ge H$ үед ирцэд орсон гэж батлах.
\end{enumerate}

Эцсийн шийдвэр:

\[
\text{present}_i =
\begin{cases}
1, & h_i \geq H \\
0, & h_i < H
\end{cases}
\]

Олон удаагийн баталгаажуулалт нь нэг удаагийн алдаанаас үүдэх FP болон FN эрсдэлийг бууруулна.

\subsection{Туршилтын орчин ба дахин давтагдах нөхцөл (Reproducibility)}
Энэхүү судалгааг дахин гүйцэтгэх боломжийг хангахын тулд ашигласан техник, орчны нөхцөл, хичээлийн цагийн зохион байгуулалтыг нэг мөр тэмдэглэв.

\begin{table}[H]
  \centering
  \caption{Туршилтын орчны үндсэн мэдээлэл}
  \label{tab:reproducibility-setup}
  \begin{tabular}{|l|p{9cm}|}
    \hline
    \textbf{Үзүүлэлт} & \textbf{Утга} \\
    \hline
    Тооцоолох төхөөрөмж & ASUS ROG Strix G513QM, AMD Ryzen 9 5900HX (16 CPUs), 16 GB RAM \\
    \hline
    Камер & Hikvision Mini PTZ camera \\
    \hline
    Сүлжээний төхөөрөмж & TP-Link TL-SF1005P switch \\
    \hline
    Оролтын зурагны нягтрал & \begin{tabular}[t]{@{}l@{}}1920\(\times\)1080 (камераас авсан эх зураг)\\Танилтын өмнө 160\(\times\)160 болгон стандартчилсан\end{tabular} \\
    \hline
    Анги танхимын хэмжээ & 6.3 м \(\times\) 6.3 м \\
    \hline
    Гэрэлтүүлгийн нөхцөл & \begin{tabular}[t]{@{}l@{}}Барилгын сүүдэр талд байрладаг\\Үдээс хойш нарны шууд тусгал нэмэгддэг\end{tabular} \\
    \hline
    Ангийн хэмжээ & 22 оюутан (\texttt{attendance/class\_roster.csv}) \\
    \hline
    Ажиглалтын өдөр тутмын хугацаа & \begin{tabular}[t]{@{}l@{}}08:00--16:25 (1 өдөр)\\1-р паар: 08:00--09:30\\2-р паар: 09:35--11:05\\3-р паар: 11:10--12:40\\4-р паар: 13:20--14:50\\5-р паар: 14:55--16:25\end{tabular} \\
    \hline
  \end{tabular}
\end{table}

\subsection{Periodic параметрийн сонголтын үндэслэл}
Periodic горимд 2--3 минут тутам нэг зураг авахад нэг хичээлийн (90 минут) хугацаанд дунджаар 30--45 зураг цугларна. Иймд ирцийг нэг кадраар биш, олон ажиглалтаар баталгаажуулахын тулд hit-д суурилсан шийдвэр хэрэглэсэн.

\begin{table}[H]
  \centering
  \caption{Periodic параметрүүд ба сонгосон үндэслэл}
  \label{tab:periodic-params-rationale}
  \begin{tabular}{|l|c|p{8.5cm}|}
    \hline
    \textbf{Параметр} & \textbf{Утга} & \textbf{Сонгосон шалтгаан} \\
    \hline
    interval (\(\Delta t\)) & 2--3 мин & CPU ачааллыг бууруулж, нэг хичээлд хангалттай олон ажиглалт (30--45 зураг) цуглуулах. \\
    \hline
    required\_hits (\(H\)) & 5--10 (зорилтот) & Нэг удаагийн буруу танилтаар ирц батлагдахаас сэргийлэх; олон ажиглалтаар баталгаажуулах. \\
    \hline
    min\_confirm\_score (\(\tau\)) & 0.75 (эхлэл утга) & Бага оноотой эргэлзээтэй танилтуудыг hit тооллоос хасч FP бууруулах. \\
    \hline
    min\_gap (\(g\)) & 2 мин & Ойрхон дараалсан ижил танилтыг давхар тоолохоос сэргийлэх. \\
    \hline
  \end{tabular}
\end{table}

Практик ажиллагаанд босго параметрүүд хоорондоо харилцан нөлөөтэй байдаг: \(\tau\)-г бууруулахад FP өсөх, \(H\)-г өсгөхөд FN өсөх хандлагатай. Иймээс параметрийн эцсийн утгыг орчны өгөгдөл дээр туршилтаар тогтоох шаардлагатай.

\subsection{Параметрийн мэдрэмжийн туршилт}
Судалгаанд параметрийн утгыг зөвхөн таамгаар тогтоохгүй, validation туршилтаар харьцуулж шалгасан. Enrollment өгөгдөл дээр leave-one-out үнэлгээ хийхэд hybrid болон k-NN арга 98.35\%--98.72\% gated accuracy үзүүлж, хамгийн тогтвортой хослол нь \(k=7\), \(\alpha=0.7\) байв. Cosine арга threshold-оос илүү мэдрэмтгий байсан бол SVM нь энэхүү жижиг өгөгдлийн нөхцөлд илүү тогтворгүй гарсан.

Энэ нь босго параметрүүдийг нэг удаагийн туршилтаар бус, бодит өгөгдлийн тархалтад тулгуурлан тохируулах шаардлагатайг харуулна. Иймээс төслийн туршилтад \(\tau\), margin, \(k\), \(\alpha\)-гийн нөлөөг тусад нь ажиглаж, хамгийн өндөр gated accuracy болон практик ажиллагааны хугацааны хослолыг сонгосон.

\section{Туршилтын аргачлал ба үнэлгээ}

\subsection{Үнэлгээний протокол}

Үнэлгээг дараах дарааллаар гүйцэтгэнэ:

\begin{enumerate}
  \item Системийн гаргасан ирцийг хадгалах.
  \item Гараар бүртгэсэн ирцийн мэдээллийг харьцуулах утга болгон бэлтгэх.
  \item Оюутан тус бүрээр TP, FP, TN, FN тооцох.
  \item Танигдаагүй тохиолдлыг тусад нь бүртгэх.
  \item Нэг өдрийн болон олон өдрийн нэгтгэл гаргах.
\end{enumerate}

\subsection{Үнэлгээний хэмжүүрүүд}

Системийн гүйцэтгэлийг Accuracy, Precision, Recall, F1-score, FAR, TAR зэрэг хэмжүүрээр үнэлнэ. Мөн гүйцэтгэлийн хугацааны үзүүлэлтүүд (нэг зураг боловсруулах хугацаа, нэг нүүр таних хугацаа, нийт мөчлөгийн хугацаа)-ийг хэмжиж, бодит хэрэглээний боломжийг үнэлнэ.

\subsection{Ёс зүй ба хувийн мэдээлэл хамгаалалт}
Сурагчдын царай нь биометр өгөгдөлд хамаарах тул зураг, эмбеддинг, ирцийн файлыг зөвхөн локал орчинд хадгалж боловсруулах зарчим баримтална. Системийг бодит ангид ашиглахын өмнө багш, сургуулийн захиргаа болон шаардлагатай тохиолдолд эцэг эх, асран хамгаалагчийн зөвшөөрөл, мэдээлэл ашиглах зорилго, хадгалалтын хугацаа, хандах эрхийг тодорхойлох шаардлагатай.

Энэхүү судалгааны хүрээнд цуглуулсан өгөгдлийг зөвхөн судалгааны зорилгоор ашиглаж, гаднын серверт автоматаар дамжуулахгүйгээр боловсруулах зарчимд тулгуурласан. Цаашдын үйлдвэрлэлийн түвшний нэвтрүүлэлтэд өгөгдөл устгах журам, access control, log retention, consent tracking, school policy compliance зэрэг нэмэлт хамгаалалтууд зайлшгүй шаардагдана.
\FloatBarrier